\section{Derivative (導數)}
\subsection{Notation}
\subsubsection{Leibniz's notation (萊布尼茲符號) for derivatives}
Suppose a dependent variable $y$ represents a function $f$ of an independent variable $x$, that is,
\[y=f(x).\]
Then, the derivative function of $y$ or $f$ (with respect to (w.r.t.) $x$) can be written as
\[\frac{\mathrm{d}y}{\mathrm{d}x},\quad\frac{\mathrm{d}}{\mathrm{d}x}y,\quad\frac{\mathrm{d}\qty(f(x))}{\mathrm{d}x},\quad\text{or\ }\frac{\mathrm{d}}{\mathrm{d}x}\qty(f(x)),\]
in which $\frac{\mathrm{d}}{\mathrm{d}x}$ is called a differential operator (微分運算子) or a derivative operator (導數運算子);

the $n$th derivative function of $y$ or $f$ (with respect to $x$) can be written as
\[\frac{\mathrm{d}^ny}{\mathrm{d}x^n},\quad\frac{\mathrm{d}^n}{\mathrm{d}x^n}y,\quad\frac{\mathrm{d}^n\qty(f(x))}{\mathrm{d}x^n},\quad\tx{or\ }\frac{\mathrm{d}^n}{\mathrm{d}x^n}\qty(f(x)),\]
in which $\frac{\mathrm{d}^n}{\mathrm{d}x^n}$ is called a differential operator or a derivative operator.
\subsubsection{Leibniz's notation for partial derivatives}
Suppose a dependent variable $y$ represents a function $f$ of an independent variable vector $\mb{x}=(x_1,x_2,\dots x_n)$, that is,
\[y=f(\mb{x}).\]
Then, the partial derivative function of $y$ or $f$ with respect to $x_i$ can be written as
\[\frac{\partial y}{\partial x_i},\quad\frac{\partial }{\partial x_i}y,\quad\frac{\partial \qty(f(x))}{\partial x_i},\quad\text{or\ }\frac{\partial }{\partial x_i}\qty(f(x)),\]
in which $\frac{\partial}{\partial x_i}$ is called a partial differential operator (偏微分運算子) or a partial derivative operator (偏導數運算子);

the $n$th partial derivative function of $y$ or $f$ with respect to $x_i$ can be written as
\[\frac{\partial^ny}{\partial x_i^{\phantom{i}n}},\quad\frac{\partial^n}{\partial x_i^{\phantom{i}n}}y,\quad\frac{\partial^n\qty(f(x))}{\partial x_i^{\phantom{i}n}}\text{or\ }\frac{\partial^n}{\partial x_i^{\phantom{i}n}}\qty(f(x)),\]
in which $\frac{\partial^n}{\partial x_i^{\phantom{i}n}}$ is called a partial differential operator or a partial derivative operator;

the $n$th mixed partial derivative function of $y$ or $f$ $m_{i_k}$ times with respect to $x_{i_k}$, and then $m_{i_{k-1}}$ times with respect to $x_{i_{k-1}}$, $\ldots$, and then $m_{i_1}$ times with respect to $x_{i_}$, in which $\sum_{j=1}^km_{i_j}=n$, can be written as
\[\frac{\partial^ny}{\partial^{m_{i_1}}x_{i_1}^{\phantom{i_1}m_{i_1}}\partial^{m_{i_2}}x_{i_2}^{\phantom{i_2}m_{i_2}}\dots \partial^{m_{i_k}}x_{i_k}^{\phantom{i_k}m_{i_k}}},\]
\[\frac{\partial^n}{\partial^{m_{i_1}}x_{i_1}^{\phantom{i_1}m_{i_1}}\partial^{m_{i_2}}x_{i_2}^{\phantom{i_2}m_{i_2}}\dots \partial^{m_{i_k}}x_{i_k}^{\phantom{i_k}m_{i_k}}}y,\]
\[\frac{\partial^nf\qty(\mb{x})}{\partial^{m_{i_1}}x_{i_1}^{\phantom{i_1}m_{i_1}}\partial^{m_{i_2}}x_{i_2}^{\phantom{i_2}m_{i_2}}\dots \partial^{m_{i_k}}x_{i_k}^{\phantom{i_k}m_{i_k}}},\quad \tx{or}\]
\[\frac{\partial^n}{\partial^{m_{i_1}}x_{i_1}^{\phantom{i_1}m_{i_1}}\partial^{m_{i_2}}x_{i_2}^{\phantom{i_2}m_{i_2}}\dots \partial^{m_{i_k}}x_{i_k}^{\phantom{i_k}m_{i_k}}}f(\mb{x}),\]
in which $\frac{\partial^n}{\partial^{m_{i_1}}x_{i_1}^{\phantom{i_1}m_{i_1}}\partial^{m_{i_2}}x_{i_2}^{\phantom{i_2}m_{i_2}}\dots \partial^{m_{i_k}}x_{i_k}^{\phantom{i_k}m_{i_k}}}$ is called a partial differential operator or a partial derivative operator.
\subsubsection{Lagrange's notation (拉格朗日符號) or Prime notation for derivatives}
Suppose a dependent variable $y$ represents a function $f$ of an independent variable $x$, that is,
\[y=f(x).\]
Then, the derivative function of $y$ or $f$ (with respect to $x$) can be written as
\[y',\quad\tx{or\ }f'(x);\]
the $n$th derivative function of $y$ or $f$ (with respect to $x$) can be written as
\[y^{(n)},\quad\tx{or\ }f^{(n)}(x),\]
in which $^{(n)}$ can be replaced with $n$ primes ($'$) for usually $n<4$ (For example, $y''$ is equivalent to $y^{(2)}$, and $f''$ is equivalent to $f^{(2)}$.).
\subsubsection{Newton's notation (牛頓符號), dot notation, flyspeck notation, or fluxions for derivatives}
Suppose a dependent variable $y$ represents a function $f$ of an independent variable $t$, that is,
\[y=f(t),\]
where $t$ usually represents time.

Then, the derivative function of $y$ or $f$ (with respect to $t$) of the function $f$ can be written as
\[\dot{y},\]
the second derivative function of $y$ or $f$ (with respect to $t$) of the function $f$ can be written as
\[\ddot{y},\]
and so on.
\subsubsection{Subscript notation for partial derivatives}
Suppose a dependent variable $y$ represents a function $f$ of an independent variable vector $\mb{x}=(x_1,x_2,\dots x_n)$, that is,
\[y=f(\mb{x}).\]
Then, the partial derivative function of $y$ or $f$ with respect to $x_i$ can be written as
\[y_{x_i},\quad y'_{x_i},\quad f_{x_i},\quad \tx{or\ }f'_{x_i};\]
the $n$th partial derivative of the function $f$ with respect to $x_i$ can be written as
\[y_{x_ix_i\dots  x_i},\quad y^{(n)}_{\pht{(n)}x_ix_i\dots  x_i},\quad f_{x_ix_i\dots  x_i},\tx{or\ }\quad f^{(n)}_{\pht{(n)}x_ix_i\dots  x_i},\]
in which the subscript are $n$ $x_i$s and $^{(n)}$ can be replaced with $n$ primes ($'$) for usually $n<4$ (For example, $y''$ is equivalent to $y^{(2)}$, and $f''$ is equivalent to $f^{(2)}$.);

the $n$th mixed partial derivative of the function $f$ $m_{i_1}$ times with respect to $x_{i_1}$, and then $m_{i_2}$ times with respect to $x_{i_2}$, $\ldots$, and then $m_{i_k}$ times with respect to $x_{i_k}$, in which $\sum_{j=1}^km_{i_j}=n$, can be written as
\[y_{x_{i_1}x_{i_1}\dots  x_{i_2}x_{i_2}x_{i_2}\dots  x_{i_2}\dots  x_{i_k}x_{i_k}\dots  x_{i_k}},\quad y^{(n)}_{\pht{(n)}x_{i_1}x_{i_1}\dots  x_{i_2}x_{i_2}x_{i_2}\dots  x_{i_2}\dots  x_{i_k}x_{i_k}\dots  x_{i_k}},\]
\[f_{x_{i_1}x_{i_1}\dots  x_{i_2}x_{i_2}x_{i_2}\dots  x_{i_2}\dots  x_{i_k}x_{i_k}\dots  x_{i_k}},\quad \tx{or\ }f^{(n)}_{\pht{(n)}x_{i_1}x_{i_1}\dots  x_{i_2}x_{i_2}x_{i_2}\dots  x_{i_2}\dots  x_{i_k}x_{i_k}\dots  x_{i_k}},\]
in which the subscript are $m_{i_j}$ $x_{i_j}$s for all $x_{i_j}$.
\subsubsection{Euler's notation (歐拉符號) for derivatives}
Suppose a dependent variable $y$ represents a function $f$ of an independent variable $x$, that is,
\[y=f(x).\]
Then, the derivative function of $y$ or $f$ (with respect to $x$) can be written as
\[Df(x),\quad(Df)(x),\quad D_xf(x),\quad (D_xf)(x),\quad \mathrm{d}f(x),\quad(\mathrm{d}f)(x),\quad \mathrm{d}_xf(x),\quad \tx{or\ }(\mathrm{d}_xf)(x),\]
in which $D$, $D_x$, $\mathrm{d}$, or $\mathrm{d}_x$ is called a (unary) differential operators or a derivative operator, and  $Df$, $D_xf$, $\mathrm{d}f$, or $\mathrm{d}_xf$ are called the differential (微分子) of $y$ or $f$ (with respect to $x$);

the $n$th derivative function of $y$ or $f$ (with respect to $x$) can be written as
\[D^nf(x),\quad(D^nf)(x),\quad D^n_{\pht{(n)}xx\dots x}f(x),\quad\qty(D^n_{\pht{(n)}xx\dots x}f)(x),\quad \mathrm{d}^nf(x),\quad(\mathrm{d}^nf)(x),\quad \mathrm{d}^n_{\pht{(n)}xx\dots x}f(x),\quad \tx{or\ }\qty(\mathrm{d}^n_{\pht{(n)}xx\dots x}f)(x),\]
in which subscript are $n$ $x$s, $D^n$, $D^n_{\pht{(n)}xx\dots x}$, $\mathrm{d}^n$, or $\mathrm{d}^n_{\pht{(n)}xx\dots x}$ is called a (unary) differential operators or a derivative operators, and $D^nf$, $D^n_{\pht{(n)}xx\dots x}f$, $\mathrm{d}^nf$, or $\mathrm{d}^n_{\pht{(n)}xx\dots x}f$ is called the differential of the $(n-1)$th derivative function of $y$ or $f$.
\subsubsection{Euler's notation for partial derivatives}
Suppose a dependent variable $y$ represents a function $f$ of an independent variable vector $\mb{x}=(x_1,x_2,\dots x_n)$, that is,
\[y=f(\mb{x}).\]
Then, the partial derivative function of $y$ or $f$ with respect to $x_i$ can be written as
\[\partial_{x_i}f(\mb{x}),\quad \qty(\partial_{x_i}f)(\mb{x}),\quad \mathrm{d}_{x_i}f(\mb{x})\tx{or\ }\qty(\mathrm{d}_{x_i}f)(\mb{x}),\quad D_{x_i}f(\mb{x})\tx{or\ }\qty(D_{x_i}f)(\mb{x}),\]
in which $\partial_{x_i}$, $\mathrm{d}_{x_i}$, or $D_{x_i}$ is called a partial differential operator or a partial derivative operator;

the $n$th partial derivative function of $y$ or $f$ with respect to $x_i$ can be written as
\[\partial_{x_ix_i\dots  x_i}f(\mb{x}),\quad\partial^{(n)}_{\pht{(n)}x_ix_i\dots  x_i}f(\mb{x}),\quad D_{x_ix_i\dots  x_i}f(\mb{x}),\quad D^{(n)}_{\pht{(n)}x_ix_i\dots  x_i}f(\mb{x}),\quad \mathrm{d}_{x_ix_i\dots  x_i}f(\mb{x}),\quad\tx{or\ }\mathrm{d}^{(n)}_{\pht{(n)}x_ix_i\dots  x_i}f(\mb{x}),\]
in which subscript are $n$ $x_i$s, $^{(n)}$ can be replaced with $n$ primes ($'$) for usually $n<4$ (For example, $\partial''$ is equivalent to $\partial^{(2)}$, and $D''$ is equivalent to $D^{(2)}$.), and $\partial_{x_ix_i\dots  x_i}$, $\partial^{(n)}_{\pht{(n)}x_ix_i\dots  x_i}$, $D_{x_ix_i\dots  x_i}$, $D^{(n)}_{\pht{(n)}x_ix_i\dots  x_i}$, $\mathrm{d}D_{x_ix_i\dots  x_i}$, or $\mathrm{d}^{(n)}_{\pht{(n)}x_ix_i\dots  x_i}$ is called a partial differential operator or a partial derivative operator;

the $n$th mixed partial derivative function of $y$ or $f$ $m_{i_1}$ times with respect to $x_{i_1}$, $m_{i_1}$ times with respect to $x_{i_2}$, $\dots $, $m_{i_k}$ times with respect to $x_{i_k}$, in which $\sum_{j=1}^km_{i_j}=n$, can be written as
C\[\partial_{x_{i_1}x_{i_1}\dots  x_{i_2}x_{i_2}x_{i_2}\dots  x_{i_2}\dots  x_{i_k}x_{i_k}\dots  x_{i_k}}f(\mb{x}),\quad\partial^{(n)}_{\pht{(n)}x_{i_1}x_{i_1}\dots  x_{i_2}x_{i_2}x_{i_2}\dots  x_{i_2}\dots  x_{i_k}x_{i_k}\dots  x_{i_k}}f(\mb{x}),\]
\[ D_{x_{i_1}x_{i_1}\dots  x_{i_2}x_{i_2}x_{i_2}\dots  x_{i_2}\dots  x_{i_k}x_{i_k}\dots  x_{i_k}}f(\mb{x}),\quad D^{(n)}_{\pht{(n)}x_{i_1}x_{i_1}\dots  x_{i_2}x_{i_2}x_{i_2}\dots  x_{i_2}\dots  x_{i_k}x_{i_k}\dots  x_{i_k}}f(\mb{x}),\]
\[ \mathrm{d}_{x_{i_1}x_{i_1}\dots  x_{i_2}x_{i_2}x_{i_2}\dots  x_{i_2}\dots  x_{i_k}x_{i_k}\dots  x_{i_k}}f(\mb{x}),\quad\tx{or\ }\mathrm{d}^{(n)}_{\pht{(n)}x_{i_1}x_{i_1}\dots  x_{i_2}x_{i_2}x_{i_2}\dots  x_{i_2}\dots  x_{i_k}x_{i_k}\dots  x_{i_k}}f(\mb{x}),\]
in which subscript are $m_{i_j}$ $x_{i_j}$s for all $x_{i_j}$, $^{(n)}$ can be replaced with $n$ primes ($'$) for usually $n<4$ (For example, $\partial''$ is equivalent to $\partial^{(2)}$, and $D''$ is equivalent to $D^{(2)}$.), and $\partial_{x_ix_i\dots  x_i}$, $\partial^{(n)}_{\pht{(n)}x_ix_i\dots  x_i}$, $D_{x_ix_i\dots  x_i}$, $D^{(n)}_{\pht{(n)}x_ix_i\dots  x_i}$, $\mathrm{d}_{x_ix_i\dots  x_i}$, and $\mathrm{d}^{(n)}_{\pht{(n)}x_ix_i\dots  x_i}$ is called a partial differential operator or a partial derivative operator.
\subsection{Ordinary Derivative of Functions with Real Domain}
Let $W$ be a topological vector space. The ordinary derivative (常導數), total derivative (全導數), or derivative (導數) of a function $f(u)\colon U\subseteq\mathbb{R}\to W$ (with respect to $u$) at $x\in U$, denoted as $f'(x)$, is defined as
\[f'(x)=\lim_{h\to 0}\frac{f(x+h)-f(x)}{h}\]
if the limit exists. If such limit exists, we say $f$ is differentiable (可微的) at $x$.

We define the (first(-order)) ordinary derivative (function) (常導（函）數) or derivative (function) (導（函）數) of $f$ (with respect to $u$) as a function $f'$ with codomain $W$ such that for any $x\in U$ at which $f$ is differentiable, $f'$ maps $x$ to the derivative of $f$ at $x$.

The derivative of the $k$th(-order) ($k\in\mathbb{N}$) derivative function of $f$ at $x\in U$ is called the $(k+1)$th(-order) derivative of $f$ (with respect to $u$) at $x\in U$. The derivative function of the $k$th(-order) ($k\in\mathbb{N}$) derivative function of $f$ is called the $k+1$th(-order) derivative function of $f$ (with respect to $u$).

If for any $n\in\mathbb{N}$, the $n$th(-order) derivative of a function $f$ at a point $x$ in its domain exists, we say $f$ is infinitely differentiable at $x$.

If $f$ is differentiable at all point in $I\subseteq U$, we say $f$ is differentiable on $I$; if $f$ is differentiable on $U$, we say $f$ is differentiable. If $f^{(n-1)}$ exists and is differentiable at all point in $I\subseteq U$, we say $f$ is $n$-times differentiable on $I$; if $f^{(n-1)}$ exists and is differentiable on $U$, we say $f$ is $n$-times differentiable.

The operation of finding the derivative or derivative function is called ordinary differentiation (常微分) or differentiation (微分).

Specifically, the $0$th(-order) derivative of $f$ is $f$ itself.

The derivative of a $f$ at $x$ represents the slope (斜率) at $x$ or the lineal element at $x$ (a miniature tangent line at $x$).
\subsection{One-sided Derivative of Functions with Real Domain}
Let $W$ be a topological vector space and $f\colon U\subseteq\mathbb{R}\to W$ be a function.

The left-hand derivative of $f$ at $x\in U$, denoted as $f'_-(x)$, is defined as
\[f'_-(x)=\lim_{h\to 0^-}\frac{f(x+h)-f(x)}{h}\]
if the limit exists.

The right-hand derivative of $f$ at $x\in U$, denoted as $f'_+(x)$, is defined as
\[f'_+(x)=\lim_{h\to 0^+}\frac{f(x+h)-f(x)}{h}\]
if the limit exists.
\subsection{Ordinary derivative of Functions with Real Vector Domain}
Let $W$ be a normed vector space. A function $f(u)\colon U\subseteq\mathbb{R}^n\to W$ with $n\in\bbN$ is differentiable at $x\in U$ if there exists a bounded linear operator $A\colon\bbR^n\to W$ such that
\[\lim_{\|h\|\to 0}\frac{\|f(x+h)-f(x)-A(h)\|_W}{\|h\|}=0.\]
If there exists such an operator $A$, it is unique, so we define the (first(-order)) ordinary derivative, total derivative, or derivative of $f$ (with respect to $u$) at $x$, denoted as $Df(x)$, as $A$.

We define the (first(-order)) ordinary derivative (function) or derivative (function) of $f$ (with respect to $u$) as a function $Df$ with codomain $B(\bbR^n,W)$, in which $B(\bbR^n,W)$ is the space of all bounded linear operators from $\bbR^n$ to $W$, such that for any $x\in U$ at which $f$ is differentiable, $Df$ maps $x$ to the derivative of $f$ at $x$.

The derivative of the $k$th(-order) ($k\in\mathbb{N}$) derivative function of $f$ at $x\in U$ is called the $k+1$th(-order) derivative of $f$ (with respect to $u$) at $x\in U$, denoted as $(D^{k+1}f)(x)$. The derivative function of the $k$th(-order) ($k\in\mathbb{N}$) derivative function of $f$ is called the $k+1$th(-order) derivative function of $f$ (with respect to $u$), denoted as $D^{k+1}f$.

If $f$ is differentiable at all point in $I\subseteq U$, we say $f$ is differentiable on $I$; if $f$ is differentiable on $U$, we say $f$ is differentiable. If $D^{(n-1)}f$ exists and is differentiable at all point in $I\subseteq U$, we say $f$ is $n$-times differentiable on $I$; if $D^{(n-1)}f$ exists and is differentiable on $U$, we say $f$ is $n$-times differentiable.

The operation of finding the derivative or derivative function is called ordinary differentiation or differentiation.

Specifically, the $0$th(-order) Fréchet derivative of $f$ is $f$ itself.
\subsection{Partial derivative of Functions with Real Vector Domain}
Let $W$ be a topological vector space and $f(\mb{x})\colon U\subseteq\mathbb{R}^n\to W$ be a function with $\mb{x}=(x_1,x_2,\ldots,x_n)$, and $X$ be the set $\{x_1,x_2,\ldots,x_n\}$.

The (first(-order)) partial derivative (偏導數) $\pdv{f}{x_i}$ of $f$ with respect to $x_i\in X$ at $u\in U$ is defined as
\bma
\pdv{f}{x_i}&=\lim_{h\to 0}\frac{f(u+h\mb{e}_i)-f(u)}{h}\\
&=\frac{\mathrm{d}}{\mathrm{d}h}f(u+h\mb{e}_i)\big\vert_{h=0}
\end{aligned},\]
in which $\mb{e}_i\in\mathbb{R}^n$ is the unit vector in the direction of $x_i$.

The partial derivative of the $k$th(-order) ($k\in\mathbb{N}$) partial derivative (function) (偏導（函）數) of $f$ with respect to $ x \in X$ with respect to $ x \in X$ at $u\in U$ is called the $k+1$th(-order) partial derivative of $f$ with respect to $ x$ at $u\in U$. The partial derivative function of the $k$th(-order) ($k\in\mathbb{N}$) partial derivative function of $f$ with respect to $ x \in X$ with respect to $ x \in X$ is called the $k+1$th(-order) partial derivative function of $f$ with respect to $ x$.

The partial derivative of the $k$th(-order) partial derivative function of $f$ with respect to $ x_1\in X$ with respect to $ x_2\in X$ at $u\in U$ is called the $(k+1)$th(-order) mixed partial derivative of $f$ with respect to $ x_1, x_1,\dots  x_1$ ($k$ times) and $ x_2$ at $u\in U$. The partial derivative function of the $k$th(-order) partial derivative function of $f$ with respect to $ x_1\in X$ with respect to $ x_2\in X$ at $u\in U$ is called the $(k+1)$th(-order) mixed partial derivative function of $f$ with respect to $ x_1, x_1,\dots  x_1$ ($k$ times) and $ x_2$.

The partial derivative of the $k$th(-order) mixed partial derivative function of $f$ with respect to $ x_1, x_2,\dots  x_k\in X$ with respect to $ x_{k+1}\in X$ at $u\in U$ is called the $(k+1)$th(-order) mixed partial derivative of $f$ with respect to $ x_1, x_2,\dots  x_k, x_{k+1}$ at $u\in U$. The partial derivative function of the $k$th(-order) mixed partial derivative function of $f$ with respect to $ x_1, x_2,\dots  x_k\in X$ with respect to $ x_{k+1}\in X$ is called the $(k+1)$th(-order) mixed partial derivative function of $f$ with respect to $ x_1, x_2,\dots  x_k, x_{k+1}$.

The operation of finding the partial derivative or partial derivative function is called partial differentiation (偏微分).

Specifically, the $0$th(-order) partial derivative of $f$ is $f$ itself.
\subsection{Fréchet derivative (弗蘭歇導數)}
\subsubsection{Fréchet derivative}
Let $V$ and $W$ be normed vector spaces and $U$ be an open subset of $V$. A function $f(u)\colon U\to W$ is Fréchet differentiable (弗蘭歇可微的) or differentiable at $x\in U$ if there exists a bounded linear operator $A\colon V\to W$ such that
\[\lim_{\|h\|_V\to 0}\frac{\|f(x+h)-f(x)-A(h)\|_W}{\|h\|_V}=0.\]
If there exists such an operator $A$, it is unique, so we define the (first(-order)) Fréchet derivative, ordinary derivative, or derivative of $f$ (with respect to $u$) at $x$, denoted as $Df(x)$, as $A$.

We define the (first(-order)) Fréchet derivative (function) (弗蘭歇導（函）數), ordinary derivative (function), or derivative (function) of $f$ (with respect to $u$) as a function $Df$ with codomain $B(V,W)$, in which $B(V,W)$ is the space of all bounded linear operators from $V$ to $W$, such that for any $x\in U$ at which $f$ is Fréchet differentiable, $Df$ maps $x$ to the Fréchet derivative of $f$ at $x$.

The Fréchet derivative of the $k$th(-order) ($k\in\mathbb{N}$) Fréchet derivative function of $f$ at $x\in U$ is called the $k+1$th(-order) Fréchet derivative of $f$ (with respect to $u$) at $x\in U$, denoted as $(D^{k+1}f)(x)$. The Fréchet derivative function of the $k$th(-order) ($k\in\mathbb{N}$) Fréchet derivative function of $f$ is called the $k+1$th(-order) Fréchet derivative function of $f$ (with respect to $u$), denoted as $D^{k+1}f$.

If $f$ is Fréchet differentiable at all point in $I\subseteq U$, we say $f$ is Fréchet differentiable on $I$; if $f$ is Fréchet differentiable on $U$, we say $f$ is Fréchet differentiable. If $D^{(n-1)}f$ exists and is Fréchet differentiable at all point in $I\subseteq U$, we say $f$ is $n$-times Fréchet differentiable on $I$; if $D^{(n-1)}f$ exists and is Fréchet differentiable on $U$, we say $f$ is $n$-times Fréchet differentiable.

The operation of finding the Fréchet derivative or Fréchet derivative function is called Fréchet differentiation (弗蘭歇微分), ordinary differentiation, or differentiation.

Specifically, the $0$th(-order) Fréchet derivative of $f$ is $f$ itself.
\subsection{Smoothness (光滑性 or 平滑性)}
\subsubsection{Continuously differentiable functions}
A function is called ($n$-times) continuously differentiable iff it is ($n$-times) differentiable and its ($n$th) derivative is continuous.
\subsubsection{Smoothness}
A function $f$ with real domain that has a derivative that is continuous on its domain is said to be of class $C^1$ or continuously differentiable, denoted as $f\in C^1$, or be a $C^1$-function.

A function $f$ that has a derivative that is continuous on a subset $I$ of its domain is said to be of class $C^1$ on $I$, of class $C^1(I)$, or continuously differentiable on $I$, denoted as $f\in C^1(I)$.

A function $f$ with real domain that has a $k$th derivative that is continuous on its domain is said to be of class $C^k$ or $k$-times continuously differentiable, denoted as $f\in C^k$, or be a $C^k$-function.

A function $f$ that has a $k$th derivative that is continuous on a subset $I$ of its domain is said to be of class $C^k$ on $I$ or of class $C^k(I)$, denoted as $f\in C^k(I)$.

Generally, the term smooth function refers to a $C^{\infty}$-function, that is a function that is of class $C^k$ for any $k\in\bbR$. However, it may also mean "sufficiently differentiable" for the problem under consideration.
\subsubsection{Piecewise smoothness of real functions}
Let $f\colon D\subseteq\mathbb{R}\to\mathbb{R}$ be a function. If there exists a family $\{[a_i,b_i]\mid i\in I\}$ of closed intervals such that:
\begin{itemize}
\item for every closed interval $V$, the set $\{i\in I\mid V\cap [a_i,b_i]\neq\varnothing\}$ is finite (some sources require instead that $n(I)$ is finite, which is stronger),
\item $D\subseteq\bigcup_{i\in I}[a_i,b_i]$,
\item function $f\big\vert_{a_i}^{b_i}\colon(a_i,b_i)\to Y$ is of class $C^k((a_i,b_i))$ for every $i\in I$, and
\item there exist finite $\lim_{x\to a_i^{\phantom{i}+}}f(x)$ and $\lim_{x\to b_i^{\phantom{i}-}}f(x)$ for every $i\in I$.
\end{itemize},
we say $f$ is of class piecewise $C^k$ (or of class piecewise $C^k\qty(\bigcup_{i\in I}[a_i,b_i])$); when $k=\infty$ or sufficiently large, we say $f$ is piecewise smooth (on $\bigcup_{i\in I}U_i$).
\subsubsection{Piecewise smoothness of functions between topological spaces}
Let $X$ and $Y$ be topological spaces, and $f\colon D\subseteq X\to Y$ be a function. If there exists a locally finite (some sources require instead finite, which is stronger) family $\{U_i\mid i\in I\}$ of closed subsets of $X$ such that:
\begin{itemize}
\item $D\subseteq\bigcup_{i\in I}U_i$, and
\item there exists a function $g_i\colon U_i\to Y$ that is of class $C^k\qty(U_i)$ such that $\forall x\in\operatorname{int}\left(U_i\right)f(x)=g_i(x)$ for every $i\in I$.
\end{itemize},
we say $f$ is of class piecewise $C^k$ (or of class piecewise $C^k\qty(\bigcup_{i\in I}U_i)$); when $k=\infty$ or sufficiently large, we say $f$ is piecewise smooth (on $\bigcup_{i\in I}U_i$).
\sssc{With compact support}
$C^k_c(U)$ for $k\in\bbN\cup\{\infty\}$ denotes the vector space of $C^k(U)$-functionals with compact support. For any compact subset $K$ of $U$, $C^k(K;U)$ denotes the vector space of $C^k(U)$-functionals with compact support on $K$.
\subsection{Gateaux Differentiation (加托微分)}
\subsubsection{Gateaux derivative (加托導數) and partial derivatives}
Let $V$ be a locally convex topological vector spaces (LCTVS), $W$ be a topological vector space, $U$ be an open subset of $V$, and $f\colon U\to W$ be a function.

The (first(-order)) Gateaux derivative $df(u;\,\psi)$ of $f$ at $u\in U$ in the direction $\psi \in V$ is defined to be
\[\begin{aligned}
df(u;\,\psi) &= \lim_{\tau\to 0}\frac{f(u+\tau \psi)-f(u)}{\tau}\\
&= \frac{\mathrm{d}}{\mathrm{d}\tau}f(u+\tau \psi)\big\vert_{\tau =0}
\end{aligned}\]
if the limit exists. If for any $\psi \in V$, the Gateaux derivative exists, then it is said that $f$ is Gateaux differentiable (加托可微的) at $u$.

We define the (first(-order)) Gateaux derivative (function) (加托導（函）數) of $f$ in the direction $\psi \in V$, denoted as $df$, as a function $df\colon U\to W$, such that for any $u\in U$ at which $f$ is Gateaux differentiable, $df$ maps $u$ to the Gateaux derivative of $f$ at $u$ in the direction $\psi \in V$.

The Gateaux derivative of the $k$th(-order) ($k\in\mathbb{N}$) Gateaux derivative function of $f$ in the direction $\psi \in V$ in the direction $\psi \in V$ at $u\in U$ is called the $k+1$th(-order) Gateaux derivative of $f$ in the direction $\psi$ at $u\in U$. The Gateaux derivative function of the $k$th(-order) ($k\in\mathbb{N}$) Gateaux derivative function of $f$ in the direction $\psi \in V$ in the direction $\psi \in V$ is called the $k+1$th(-order) Gateaux derivative function of $f$ in the direction $\psi$.

The Gateaux derivative of the $k$th(-order) Gateaux derivative function of $f$ in the direction $\psi_1\in V$ in the direction $\psi_2\in V$ at $u\in U$ is called the $(k+1)$th(-order) mixed Gateaux derivative of $f$ in the direction $\psi_1,\psi_1,\dots \psi_1$ ($k$ times) and $\psi_2$ at $u\in U$. The Gateaux derivative function of the $k$th(-order) Gateaux derivative function of $f$ in the direction $\psi_1\in V$ in the direction $\psi_2\in V$ is called the $(k+1)$th(-order) mixed Gateaux derivative function of $f$ in the direction $\psi_1,\psi_1,\dots \psi_1$ ($k$ times) and $\psi_2$.

The Gateaux derivative of the $k$th(-order) mixed Gateaux derivative function of $f$ in the direction $\psi_1,\psi_2,\dots \psi_k\in V$ in the direction $\psi_{k+1}\in V$ at $u\in U$ is called the $(k+1)$th(-order) mixed Gateaux derivative of $f$ in the direction $\psi_1,\psi_2,\dots \psi_k,\psi_{k+1}$ at $u\in U$. The Gateaux derivative function of the $k$th(-order) mixed Gateaux derivative function of $f$ in the direction $\psi_1,\psi_2,\dots \psi_k\in V$ in the direction $\psi_{k+1}\in V$ is called the $(k+1)$th(-order) mixed Gateaux derivative function of $f$ in the direction $\psi_1,\psi_2,\dots \psi_k,\psi_{k+1}$.

The operation of finding the Gateaux derivative or Gateaux derivative function is called Gateaux differentiation (加托微分).

Specifically, the $0$th(-order) Gateaux derivative of $f$ is $f$ itself.
\subsubsection{Partial derivative}
Let $V$ be a locally convex topological vector spaces (LCTVS), $W$ be a topological vector space, $U$ be an open subset of $V$, and $f\colon U\to W$ be a function, $\mb{x}$ be the independent variable vector of $f$, and $X$ be the set of all independent variables of $f$.

The (first(-order)) partial derivative of $f$ with respect to $x_i\in X$ at $u\in U$ is defined as the Gateaux derivative of $f$ in the direction of $x_i$ at $u$ if it exists.

The partial derivative of the $k$th(-order) ($k\in\mathbb{N}$) partial derivative (function) of $f$ with respect to $ x \in X$ with respect to $ x \in X$ at $u\in U$ is called the $k+1$th(-order) partial derivative of $f$ with respect to $ x$ at $u\in U$. The partial derivative function of the $k$th(-order) ($k\in\mathbb{N}$) partial derivative function of $f$ with respect to $ x \in X$ with respect to $ x \in X$ is called the $k+1$th(-order) partial derivative function of $f$ with respect to $ x$.

The partial derivative of the $k$th(-order) partial derivative function of $f$ with respect to $ x_1\in X$ with respect to $ x_2\in X$ at $u\in U$ is called the $(k+1)$th(-order) mixed partial derivative of $f$ with respect to $ x_1, x_1,\dots  x_1$ ($k$ times) and $ x_2$ at $u\in U$. The partial derivative function of the $k$th(-order) partial derivative function of $f$ with respect to $ x_1\in X$ with respect to $ x_2\in X$ at $u\in U$ is called the $(k+1)$th(-order) mixed partial derivative function of $f$ with respect to $ x_1, x_1,\dots  x_1$ ($k$ times) and $ x_2$.

The partial derivative of the $k$th(-order) mixed partial derivative function of $f$ with respect to $ x_1, x_2,\dots  x_k\in X$ with respect to $ x_{k+1}\in X$ at $u\in U$ is called the $(k+1)$th(-order) mixed partial derivative of $f$ with respect to $ x_1, x_2,\dots  x_k, x_{k+1}$ at $u\in U$. The partial derivative function of the $k$th(-order) mixed partial derivative function of $f$ with respect to $ x_1, x_2,\dots  x_k\in X$ with respect to $ x_{k+1}\in X$ is called the $(k+1)$th(-order) mixed partial derivative function of $f$ with respect to $ x_1, x_2,\dots  x_k, x_{k+1}$.

The operation of finding the partial derivative or partial derivative function is called partial differentiation.

Specifically, the $0$th(-order) partial derivative of $f$ is $f$ itself.
\subsection{Applicational interpretation}
\subsubsection{Rate of change}
If $y=f(x)$, then $f'(a)$ where defined is called the instantaneous rate of change of $y$ with repsect to $x$ at $a$.
A function $f\colon U\to W$ is Fréchet differentiable (弗蘭歇可微的) or differentiable at $x\in U$ if there exists a bounded linear operator $A\colon V\to W$ such that
\[\lim_{\|h\|_V\to 0}\frac{\|f(x+h)-f(x)-A(h)\|_W}{\|h\|_V}=0.\]
If there exists such an operator $A$, it is unique, so we define the (first(-order)) Fréchet derivative, ordinary derivative, or derivative of $f$ (with respect to $u$) at $x$, denoted as $Df(x)$, as $A$.

We define the (first(-order)) Fréchet derivative (function) (弗蘭歇導（函）數), ordinary derivative (function), or derivative (function) of $f$ (with respect to $u$) as a function $Df$ with codomain $B(V,W)$, in which $B(V,W)$ is the space of all bounded linear operators from $V$ to $W$, such that for any $x\in U$ at which $f$ is Fréchet differentiable, $Df$ maps $x$ to the Fréchet derivative of $f$ at $x$.

The Fréchet derivative of the $k$th(-order) ($k\in\mathbb{N}$) Fréchet derivative function of $f$ at $x\in U$ is called the $k+1$th(-order) Fréchet derivative of $f$ (with respect to $u$) at $x\in U$, denoted as $(D^{k+1}f)(x)$. The Fréchet derivative function of the $k$th(-order) ($k\in\mathbb{N}$) Fréchet derivative function of $f$ is called the $k+1$th(-order) Fréchet derivative function of $f$ (with respect to $u$), denoted as $D^{k+1}f$.

If $f$ is Fréchet differentiable at all point in $I\subseteq U$, we say $f$ is Fréchet differentiable on $I$; if $f$ is Fréchet differentiable on $U$, we say $f$ is Fréchet differentiable. If $D^{(n-1)}f$ exists and is Fréchet differentiable at all point in $I\subseteq U$, we say $f$ is $n$-times Fréchet differentiable on $I$; if $D^{(n-1)}f$ exists and is Fréchet differentiable on $U$, we say $f$ is $n$-times Fréchet differentiable.

The operation of finding the Fréchet derivative or Fréchet derivative function is called Fréchet differentiation (弗蘭歇微分), ordinary differentiation, or differentiation.

Specifically, the $0$th(-order) Fréchet derivative of $f$ is $f$ itself.

If $y=f(x)$ is defined on some interval $[a,b]$ with $b>a$, then $\frac{f(b)-f(a)}{b-a}$ is called the average rate of change of $y$ with repsect to $x$ over $[a,b]$.
\subsubsection{Physical interpretation}
Let $f$ be a function with domain being $\bbR$ or $\bbR_{\geq 0}$. If we interprets $f(t)$ as the position of $y$ at time $t$, then:
\bit
\item $f$ is called the position function of $y$,
\item if $f(x)$ is defined on some interval $[a,b]$ with $b>a$, then $f(b)-f(a)$ is called the displacement of $y$ over $[a,b]$,
\item $f'$ is called the velocity function of $y$,
\item $f'(a)$ where defined is called the instantaneous velocity of $y$ at $a$,
\item if $f(x)$ is defined on some interval $[a,b]$ with $b>a$, then $\frac{f(b)-f(a)}{b-a}$ is called the average velocity of $y$ over $[a,b]$,
\item $f''$ is called the acceleration function of $y$,
\item $f''(a)$ where defined is called the instantaneous acceleration of $y$ at $a$,
\item if $f'(x)$ is defined on some interval $[a,b]$ with $b>a$, then $\frac{f'(b)-f'(a)}{b-a}$ is called the average acceleration of $y$ over $[a,b]$,
\item $f'''$ is called the jerk function of $y$,
\item $f'''(a)$ where defined is called the instantaneous jerk of $y$ at $a$, and
\item if $f''(x)$ is defined on some interval $[a,b]$ with $b>a$, then $\frac{f''(b)-f''(a)}{b-a}$ is called the average jerk of $y$ over $[a,b]$.
\eit
\subsubsection{Biological interpretation}
Let $f$ be a function with domain being $\bbR$ or $\bbR_{\geq 0}$. If we interprets $f(t)$ as the number of some population $y$ at time $t$, then:
\bit
\item $f$ is called the growth (function) of $y$,
\item $f'(a)$ where defined is called the instantaneous rate of growth of $y$ at $a$, and
\item if $f(x)$ is defined on some interval $[a,b]$ with $b>a$, then $\frac{f(b)-f(a)}{b-a}$ is called the average rate of growth of $y$ over $[a,b]$.
\eit
\subsubsection{Economical interpretation}
Let $f$ be a function with domain being $\bbR$ or $\bbR_{\geq 0}$. If we interprets $f(t)$ as the benefit (aka revenue) or cost of consuming or producing $y$ of quantity $t$, then:
\bit
\item $f$ is called benefit (aka revenue) or cost (function) of $y$, and
\item $f'(a)$ where defined is called the margin benefit (aka revenue) or cost of $y$ at $a$.
\eit
\subsection{Differential}
For a normed vector-valued function $f$ of an independent variable $x$, we define the differential of it, $\dd{f}$, in terms of the differential of $x$, denoted as $\dd{x}$, as
\[\dd{f}=f'\dd{x}.\]
For a normed vector-valued function $f$ of independent variables $x_1,x_2,\ldots,x_n$, we define the (total) differential of it, $\dd{f}$, in terms of the differentials of $x_1,x_2,\ldots,x_n$, denoted as $\dd{x_1},\dd{x_2},\ldots,\dd{x_n}$, as
\[\dd{f}=\sum_{k=1}^nf_{x_k}\dd{x_k}.\]
\subsection{Antiderivative (反導函數), inverse derivative, primitive function, primitive integral, or indefinite integral (不定積分)}
\subsubsection{Definition}
An antiderivative (aka inverse derivative, primitive function, primitive integral, or indefinite integral) of a function $f(x)$ on subset $I$ of the domain of $f$ is a differentiable function $F(x)$ such that
\[F'(x)=f(x)\]
for all $x\in I$.

The process of solving for antiderivatives is called antidifferentiation (反微分) or indefinite integration (不定積分).
\subsubsection{Theorem}
If a function $F(x)$ with codomain $Y$ is an antiderivative of a function $f(x)$ on subset $I$ of the domain of $f$, then the family of all antiderivatives of $f(x)$ on $I$ is
\[\{F(x)+C\mid C\in Y\}.\]
\subsubsection{Notation}
If a function $F(x)$ with codomain $Y$ is an antiderivative of a function $f(x)$ on subset $I$ of the domain of $f$, we write
\[\int f(x)\dd{x}=F(x)+C,\]
where $C$ denotes an arbitrary constant in $Y$ that does not depend on $x$, called constant of integration. The symbol $\int$ is called integral sign, introduced by Leibniz.

An indefinite integral of the form
\[\int f(x)\dd{x}=
\begin{cases}
F(x)+C_1,\quad & x\in A_1\\
F(x)+C_2,\quad & x\in A_2\\
\vdots
F(x)+C_n,\quad & x\in A_n\\
(\ldots)
\end{cases}\]
with an interval $A$ such that $\bigcup_{i=1}^nA_i\subseteq A$ and that $A\setminus\bigcup_{i=1}^nA_i$ is a set $B$ of measure zero and that $F(x)$ is not defined on $C$ and defined on $\bigcup_{i=1}^nA_i$, is usually abbreviated as
\[\int f(x)\dd{x}=
\begin{cases}
F(x)+C,\quad & x\in A\\
\ldots
\end{cases}\]
where $C$ is to be understood as notation for a locally constant function of $x$.

When computing indefinite integral, we usually absorb constants into $C$ without changing its symbol, that is,
\[\int f(x)\dd{x}=\ldots=F(x)+a+C=F(x)+C.\]
\subsection{Taylor series}
\subsubsection{Taylor series (泰勒級數) or Taylor expansion (泰勒展開)}
Let $f\colon D\subseteq\mathbb{R}\to\mathbb{R}$ be a function and $I=(a-R,a+R)\subseteq D$ be an interval such that $f\in C^{\infty}(I)$, then the Taylor series of $f$ at $a\in I$ (or about $a$, near $a$, or centered at $a$) is
\[\sum_{n=0}^{\infty}\frac{f^{(n)}(a)}{n!}(x-a)^n,\]
of which convergence is not guaranteed, and convergence to $f$ is not guaranteed even if it's convergent.

Taylor series of at $0$ is called Maclaurin series (馬克勞林級數) or Maclaurin expansion (馬克勞林展開).
\subsubsection{Taylor polynomial}
Let $f\colon D\subseteq\mathbb{R}\to\mathbb{R}$ be a function and $I=(a-R,a+R)\subseteq D$ be an interval where $f$ is $k$ times differentiable, then the $k$-th degree (or order) Taylor polynomial (or approximation) of $f$ at $a$ (or about $a$, near $a$, or centered at $a$) is
\[T_n(x)=\sum_{n=0}^k\frac{f^{(n)}(a)}{n!}(x-a)^n,\]
in which the first degree one is also called linear approximation, linearization, or tangent line (approximation), and
\[R_n(x)=f(x)-T_n(x)\]
is called remainder.

Taylor polynomial of at $0$ is called Maclaurin polynomial.
\subsubsection{Real analyticity (實解析性)}
A real-valued function $f$ is real analytic at a point $x_0$ in its domain, if it is infinitely differentiable at $x_0$ and that the Taylor expansion of $f$ at $x_0$ converges to $f(x)$ pointwise for any $x$ in a neighborhood of $x_0$.

A real-valued function is called to be real analytic on an open set $I$ that is a subset of its domain if it is real analytic at any point in $I$.

A real-valued function is called to be real analytic if it is real analytic at any point in its domain.
\subsection{Differential equation (微分方程)}
\subsubsection{Differential Equation}
An equation containing the derivatives of one or more unknown functions or dependent variables, with respect to one or more independent variables, is said to be a differential equation (DE).
\subsubsection{Ordinary differential equation (ODE) (常微分方程)}
If a differential equation contains only ordinary derivatives of one or more unknown functions with respect to a single independent variable, it is said to be an ordinary differential equation (ODE).
\subsubsection{Partial differential equation (PDE) (偏微分方程)}
If a DE contains partial derivatives of one or more unknown functions of two or more independent variables is called a partial differential equation (PDE).
\subsubsection{Order of a DE}
The order of a differential equation is the order of the highest derivative in the equation.
\subsection{Graph-related}
\subsubsection{Critical point (臨界點)}
Let $f$ be a function with domain being a subset of a LCTVS and codomain $\bbR$ and \( c \) be a point in the domain of $f$, if \( f'(c) = 0 \) or \( f' \) does not exist at \( c \), then \( c \) is a critical point, of \( f \), and $f(c)$ is called a critical value of $f$.

If the domain of $f$ is a subset of $\bbR$ or $\bbC$, a critical point is also called a critical number.
\subsubsection{Relative extremum (相對極值), local extremum (局部極值), or extremum (極值)}
Let $f\colon D\subseteq X\to\bbR$ be a function with $X$ being a topological space. It is said that a relative maximum (相對極大值) or maximum (極大值) \( f(c) \) of \(f\) occurs at $c\in D$ if there exists an open subset $I\ni c$ of $X$ such that \( \forall x\in I\cap D\colon f(c) \geq f(x) \).

Let $f\colon D\subseteq X\to\bbR$ be a function with $X$ being a topological space. It is said that a relative minimum (相對極小值) or minimum (極小值) \( f(c) \) of \(f\) occurs at $c\in D$ if there exists an open subset $I\ni c$ of $X$ such that \( \forall x\in I\cap D\colon f(c) \leq f(x) \).

Relative maximum and relative minimum are collectively called relative extreme.

Let $f\colon D\subseteq X\to\bbR$ be a function with $X$ being a topological space. It is said that a strict relative maximum or strict maximum \( f(c) \) of \(f\) occurs at $c\in D$ if there exists an open subset $I\ni c$ of $X$ such that \( \forall x\in I\cap D\colon f(c) > f(x) \).

Let $f\colon D\subseteq X\to\bbR$ be a function with $X$ being a topological space. It is said that a strict relative minimum or strict minimum \( f(c) \) of \(f\) occurs at $c\in D$ if there exists an open subset $I\ni c$ of $X$ such that \( \forall x\in I\cap D\colon f(c) < f(x) \).

Strict relative maximum and strict relative minimum are collectively called strict relative extreme.
\subsubsection{Absolute extremum (絕對極值 or 最值) or global extremum (全域極值)}
Let $f\colon D\to\bbR$ be a function. It is said that a absolute maximum (絕對極大值 or 最大值) \( f(c) \) of \(f\) occurs at $c\in D$ if \( \forall x\in D\colon f(c) \geq f(x) \).

Let $f\colon D\to\bbR$ be a function. It is said that a absolute minimum (絕對極小值 or 最小值) \( f(c) \) of \(f\) occurs at $c\in D$ if \( \forall x\in D\colon f(c) \leq f(x) \).

Absolute maximum and absolute minimum are collectively called absolute extreme.
\subsubsection{Saddle point (鞍點) or minimax}
Let $f$ be a function with codomain $\bbR$ and \( c \) be a point in the domain of $f$ , if \( f'(c) = 0 \) and $c$ is not a local extremum of $f$, then $c$ is a saddle point of $f$.
\subsubsection{Concavity (凹性)}
Let $V$ be a vector space and function $f\colon J\subseteq V\to\bbR$ be continuous on a convex set $I\subseteq J$.
\bit
\item If for any $x,y\in I$ and $\alpha\in(0,1)$, $f((1-\alpha)x+\alpha y)<(1-\alpha)f(x)+\alpha f(y)$, often called "$f$ lies strictly below all of its secant lines on $I$", then $f$ is called to be strictly concave up/upward/upwards (CU), strictly convex, strictly convex up/upward/upwards, or strictly cup-shaped.
\item If for any $x,y\in I$ and $\alpha\in(0,1)$, $f((1-\alpha)x+\alpha y)>(1-\alpha)f(x)+\alpha f(y)$, often called "$f$ lies strictly above all its secant lines on $I$", then $f$ is called to be strictly concave, strictly concave down/downward/downwards (CD), strictly convex down/downward/downwards, or strictly cap-shaped.
\item If for any $x,y\in I$ and $\alpha\in(0,1)$, $f((1-\alpha)x+\alpha y)\leq (1-\alpha)f(x)+\alpha f(y)$, often called "$f$ lies weakly below all of its secant lines on $I$", then $f$ is called to be weakly concave up/upward/upwards (CU), weakly convex, weakly convex up/upward/upwards, or weakly cup-shaped.
\item If for any $x,y\in I$ and $\alpha\in(0,1)$, $f((1-\alpha)x+\alpha y)\geq (1-\alpha)f(x)+\alpha f(y)$, often called "$f$ lies weakly above all of its secant lines on $I$", then $f$ is called to be weakly concave, weakly concave down/downward/downwards (CD), weakly convex down/downward/downwards, or weakly cap-shaped.
\eit
The terms without weakly and strictly are define to be either weakly or strictly in different contexts; here we define them to be strictly.
\subsubsection{Point of inflection or inflection point (反曲點 or 拐點)}
Let $V$ be a vector space, $f\colon J\subseteq V\to\bbR$ be a function, and $I\subseteq J$ be a convex set. If $f$ is concave upward on $(a,b)$ and concave downward on $(b,c)$, or, concave downward on $(a,b)$ and concave upward on $(b,c)$, then $(b,f(b))$ is called an inflection point of $f$.
\subsubsection{Kink}
A kink is a point $x$ on the graph of a function $f$ such that $f$ is continuous but not differentiable at $x$.
\subsubsection{Corner or sharp corner}
A corner or sharp corner is a point $x$ on the graph of a real function $f$ such that $f$ is continuous but not differentiable at $x$ and that the left-hand derivative and the right-hand derivative of $f$ at $x$ both exist and are finite.
\subsubsection{Vertical tangent}
A vertical tangent is a point $x$ on the graph of a real function $f$ such that $f$ is continuous at $x$ and that the left-hand derivative and the right-hand derivative of $f$ at $x$ are both $\infty$ or are both $-\infty$.
\subsubsection{Cusp}
A cusp is a point $x$ on the graph of a real function $f$ such that:
\bit
\item $f$ is continuous at $x$,
\item at least one of the left-hand derivative and the right-hand derivative of $f$ at $x$ are $\infty$ or $-\infty$,
\item $x$ is not a vertical tangent of $f$, and
\item the left-hand derivative and the right-hand derivative of $f$ at $x$ are both either finite, $\infty$, or $-\infty$.
\eit